% \documentclass[conference]{IEEEtran}
% \usepackage{cite}
% \usepackage{amsmath,amssymb,amsfonts}
% \usepackage{algorithmic}
% \usepackage{algorithm}
% \usepackage{graphicx}
% \usepackage{textcomp}
% \usepackage{xcolor}
% \usepackage{booktabs}

% \begin{document}

% \title{Parallel Project Report: Parallel and Distributed Hierarchical Clustering for Large Datasets}

% \author{\IEEEauthorblockN{Antonio Possemato (ID: 244262)}
% antonio.possemato@studenti.unitn.it
% \\Course: Introduction to Parallel Computing (2025–2026)}
\documentclass[conference]{IEEEtran}
\usepackage{cite}
\usepackage{amsmath,amssymb,amsfonts}
\usepackage{algorithmic}
\usepackage{algorithm}
\usepackage{graphicx}
\usepackage{textcomp}
\usepackage{xcolor}
\usepackage{booktabs}

\begin{document}

\title{Parallel Project Report: Parallel and Distributed Hierarchical Clustering for Large Datasets}

\author{\IEEEauthorblockN{Antonio Possemato (ID: 244262)}
antonio.possemato@studenti.unitn.it
\\Course: Introduction to Parallel Computing (2025–2026)}

\maketitle

\begin{abstract}
This project presents a parallel implementation of Single Linkage Hierarchical Clustering optimized for distributed multicore architectures. We employ a hybrid MPI and OpenMP approach to overcome the $O(N^2)$ memory and computational bottlenecks of standard sequential algorithms. The local shared-memory phase leverages the SHRINK algorithm, dividing the dataset into overlapping subsets and merging partial dendrograms using a Union-Find data structure. The distributed MPI phase manages global data partitioning and communication via a binomial tree reduction. We evaluate our approach on the SUSY dataset using a distributed cluster. Results demonstrate the effectiveness of our hybrid partitioning strategy in balancing computational load and minimizing memory footprint. 
\end{abstract}

\begin{IEEEkeywords}
Parallel Computing, MPI, OpenMP, Hierarchical Clustering, Single Linkage, SHRINK
\end{IEEEkeywords}

\section{Introduction}
Hierarchical clustering provides detailed structural insights into datasets without requiring a predefined number of clusters. However, standard algorithms for Single Linkage, such as SLINK or Minimum Spanning Tree (MST) based approaches, require $O(N^2)$ time and up to $O(N^2)$ space. This limits their applicability on large-scale datasets like SUSY.

This project aims to:
1) Implement a highly optimized, hybrid MPI and OpenMP Single Linkage clustering kernel.
2) Develop a multi-level data partitioning strategy to minimize memory overhead on individual nodes.
3) Analyze performance characteristics, including strong and weak scaling, across a distributed computing cluster.

\section{State of the Art}
Standard sequential approaches for Single Linkage, such as Lance-Williams or SLINK, suffer from steep memory or computational requirements. Shared-memory parallelizations often rely on Parallel Bor\r{u}vka or algorithms like SHRINK, which effectively distribute distance calculations across threads but are limited by single-node memory capacity. Existing distributed-memory implementations frequently rely on master-worker topologies or explicit global distance matrices, introducing severe network and memory bottlenecks. Our approach mitigates these limits by avoiding the explicit $O(N^2)$ matrix construction and replacing flat communication topologies with a scalable binomial tree reduction.

\section{Contribution and Methodology}
We adopt a divide-and-conquer strategy inspired by the SHRINK algorithm [1], extended to a distributed memory environment via MPI. 

\subsection{Algorithm Overview}
Our implementation decomposes the global problem into local independent clustering tasks, followed by a hierarchical merge. Algorithm 1 illustrates the hybrid strategy. The dataset $D$ is distributed among $P$ MPI ranks. Each rank processes its local block using OpenMP to produce a partial dendrogram. Finally, MPI processes communicate their partial dendrograms to build the global tree.

\begin{algorithm}
\caption{Hybrid MPI+OpenMP Hierarchical Clustering}
\begin{algorithmic}[1]
\REQUIRE Dataset $D$ of size $N$, $P$ MPI ranks, $T$ OpenMP threads per rank
\STATE \textbf{Phase 1: Block 1D Partitioning (MPI)}
\STATE Rank $p$ loads contiguous block $D_p$ of size $B = N/P$
\STATE \textbf{Phase 2: Local Clustering (OpenMP SHRINK)}
\STATE Partition $D_p$ into $K$ overlapping subsets
\STATE \textbf{\#pragma omp parallel for schedule(dynamic)}
\STATE \textbf{for} each pair of subsets $(i, j)$ \textbf{do}
\STATE \quad Compute partial dendrogram on $D_{p,i} \cup D_{p,j}$
\STATE \textbf{end for}
\STATE Merge local partial dendrograms using Union-Find to produce $T_p$
\STATE \textbf{Phase 3: Global Merge (MPI Binomial Tree)}
\STATE \textbf{for} $step = 0$ \textbf{to} $\lceil \log_2 P \rceil - 1$ \textbf{do}
\STATE \quad \textbf{if} $p \pmod{2^{step+1}} == 0$ \textbf{then}
\STATE \quad \quad \textbf{if} $p + 2^{step} < P$ \textbf{then}
\STATE \quad \quad \quad \texttt{MPI\_Recv}($T_{recv}$ from rank $p + 2^{step}$)
\STATE \quad \quad \quad Merge $T_p$ and $T_{recv}$ using Union-Find
\STATE \quad \quad \textbf{end if}
\STATE \quad \textbf{else}
\STATE \quad \quad \texttt{MPI\_Send}($T_p$ to rank $p - 2^{step}$)
\STATE \quad \quad \textbf{break} 
\STATE \quad \textbf{end if}
\STATE \textbf{end for}
\RETURN Global dendrogram on Rank 0
\end{algorithmic}
\end{algorithm}

\subsection{Data Partitioning and Memory Strategy}
We implement a static 1D block partitioning strategy. The global dataset $D$ of size $N$ is divided equally among $P$ MPI processes.
\begin{itemize}
    \item \textbf{Memory Isolation:} Each process $p$ is allocated a disjoint data block $D_p$ of size $N/P$. This bounds the memory footprint per node to $O(N/P)$, preventing Out-Of-Memory (OOM) errors on large datasets.
    \item \textbf{Zero-Communication Initialization:} Processes load their designated chunk directly from the storage file system, avoiding a root-scatter bottleneck.
\end{itemize}

\subsection{Shared-Memory Parallelization (OpenMP)}
Within each MPI process, $D_p$ is processed using OpenMP. Following the SHRINK methodology, the data is partitioned into $K$ overlapping subsets.
\begin{itemize}
    \item \textbf{Dynamic Scheduling:} We explicitly use \texttt{\#pragma omp parallel for schedule(dynamic)}. The computational cost of clustering a subset pair $(i, j)$ varies depending on the local density of the points. Dynamic scheduling prevents load imbalance among threads.
    \item \textbf{Thread-Safety:} Partial dendrograms are merged concurrently using a thread-safe Union-Find data structure.
\end{itemize}

\subsection{Distributed-Memory Merge (MPI)}
To resolve the clustering tree globally, we employ a Binomial Tree Reduction.
\begin{itemize}
    \item \textbf{Logarithmic Scalability:} A centralized master-worker gather would create an $O(P)$ network and compute bottleneck. The binomial tree bounds the communication and merge overhead to $\lceil \log_2 P \rceil$ steps.
    \item \textbf{Communication Protocol:} We transfer dendrogram edges using an Array of Structs (AoS) format. A custom MPI datatype (\texttt{MPI\_Type\_create\_struct}) encapsulates the edge properties (two integer IDs and one float distance). To handle variable-sized partial dendrograms dynamically, receiving nodes use \texttt{MPI\_Probe} and \texttt{MPI\_Get\_count} to allocate exact memory buffers, avoiding two-phase communication overhead.
    \item \textbf{Union-Find Aggregation:} At each reduction step, receiving processes merge their current dendrogram $T_p$ with the incoming $T_{recv}$ using Union-Find, discarding redundant edges without recalculating distances.
\end{itemize}

\section{Experiments and System Description}
PLACEHOLDER

\section{Results and Discussion}
PLACEHOLDER

\section{Conclusions}
PLACEHOLDER

\section{References}
[1] W. Hendrix, M. M. A. Patwary, A. Agrawal, W. Liao, and A. Choudhary, "Parallel Hierarchical Clustering on Shared Memory Platforms," ...

\end{document}